% GENERATED FILE. Edit canonical lessons, then run npm run generate:books.
% canonical-source-hash: fnv1a64:ccec88e945b6dd02
% canonical-lessons: ZH-C331-yaoguai, ZH-C331-shenxian, ZH-C331-mofa, ZH-C331-baowu, ZH-C331-cang

\chapter{Monsters and Treasure}
\label{ch:zh-monsters-and-treasure}
\hlchaptermodality{331}

\begin{chapteropening}
\textbf{By the end of this chapter:} \emph{I can talk about monsters, immortals, magic and treasure, and say people hid.}

Complete the last lesson of chapter 331: I can talk about monsters, immortals, magic and treasure, and say people hid.
\end{chapteropening}

\section[yāoguài]{yāoguài — monster; demon}

\label{lesson:ZH-C331-yaoguai}



\begin{quote}
Before the new one: say the Mandarin for a king's wife, then the Mandarin for a kingdom ruled by a king.
\end{quote}

\subsection*{You'll want to know: yāoguài}
\emph{yāoguài} — \textquotedblleft{}monster; demon\textquotedblright{}.

\begin{quote}
\emph{chuánshuō shān lǐ yǒu yí ge yāoguài} — Legend says there's a monster in the mountains.
\end{quote}

Its characters come later; for now it is a word for the ear.

\begin{cousinweb}
\emph{yāoguài} is built from \emph{yāo} and \emph{guài} 'strange', both inherited from Middle Chinese, \emph{'jew} and \emph{kweajH} (Guangyun, in Baxter's transcription), which glosses the first 'bewitching' and the second 'strange'.
\end{cousinweb}

\begin{grammarlens}[title={What you've built}]
One more word for telling what happened.
\end{grammarlens}

\subsection*{Your turn}
\begin{itemize}
\raggedright
  \item \emph{Say it:} \emph{yāoguài}
  \item \emph{Say it:} \emph{yāoguài}, once more
  \item \emph{Say it:} \emph{chuánshuō shān lǐ yǒu yí ge yāoguài} — Legend says there's a monster in the mountains.
  \item \emph{Recall:} say the Mandarin for a king's wife, then the Mandarin for a kingdom ruled by a king, then say \emph{yāoguài} again
\end{itemize}

\begin{tcolorbox}[breakable,colback=teal!4,colframe=teal!35!black,title={Before you move on}]
What does yāoguài mean? (\textquotedblleft{}Monster; demon\textquotedblright{}.) Say it once more.
\end{tcolorbox}
\section[shénxiān]{shénxiān — immortal; supernatural being}

\label{lesson:ZH-C331-shenxian}



\begin{quote}
Before the new one: say the Mandarin for a kingdom ruled by a king, then the Mandarin for a monster in a story.
\end{quote}

\subsection*{You'll want to know: shénxiān}
\emph{shénxiān} — \textquotedblleft{}immortal; supernatural being\textquotedblright{}.

\begin{quote}
\emph{zuìhòu yí ge shénxiān bāng le tā} — In the end an immortal helped him.
\end{quote}

Its characters come later; for now it is a word for the ear.

\begin{cousinweb}
\emph{shénxiān} is built from \emph{shén} 'a god' and \emph{xiān} 'an immortal', both inherited from Middle Chinese, \emph{zyin} and \emph{sjen} (Guangyun, in Baxter's transcription), which glosses the first 'a spirit' and says of the second that one who grows old and does not die is called so.
\end{cousinweb}

\begin{grammarlens}[title={What you've built}]
One more word for telling what happened.
\end{grammarlens}

\subsection*{Your turn}
\begin{itemize}
\raggedright
  \item \emph{Say it:} \emph{shénxiān}
  \item \emph{Say it:} \emph{shénxiān}, once more
  \item \emph{Say it:} \emph{zuìhòu yí ge shénxiān bāng le tā} — In the end an immortal helped him.
  \item \emph{Recall:} say the Mandarin for a kingdom ruled by a king, then the Mandarin for a monster in a story, then say \emph{shénxiān} again
\end{itemize}

\begin{tcolorbox}[breakable,colback=teal!4,colframe=teal!35!black,title={Before you move on}]
What does shénxiān mean? (\textquotedblleft{}Immortal; supernatural being\textquotedblright{}.) Say it once more.
\end{tcolorbox}
\section[mófǎ]{mófǎ — magic}

\label{lesson:ZH-C331-mofa}



\begin{quote}
Before the new one: say the Mandarin for a monster in a story, then the Mandarin for a Daoist immortal.
\end{quote}

\subsection*{You'll want to know: mófǎ}
\emph{mófǎ} — \textquotedblleft{}magic\textquotedblright{}.

\begin{quote}
\emph{nàge shénxiān huì mófǎ} — That immortal can do magic.
\end{quote}

Its characters come later; for now it is a word for the ear.

\begin{cousinweb}
\emph{mófǎ} is built from \emph{mó} 'a demon' and \emph{fǎ} 'a method', both inherited from Middle Chinese, \emph{ma} and \emph{pjop} (Guangyun, in Baxter's transcription), which glosses the first 'a demon' and the second 'a rule' and 'a norm'.
\end{cousinweb}

\begin{grammarlens}[title={What you've built}]
One more word for telling what happened.
\end{grammarlens}

\subsection*{Your turn}
\begin{itemize}
\raggedright
  \item \emph{Say it:} \emph{mófǎ}
  \item \emph{Say it:} \emph{mófǎ}, once more
  \item \emph{Say it:} \emph{nàge shénxiān huì mófǎ} — That immortal can do magic.
  \item \emph{Recall:} say the Mandarin for a monster in a story, then the Mandarin for a Daoist immortal, then say \emph{mófǎ} again
\end{itemize}

\begin{tcolorbox}[breakable,colback=teal!4,colframe=teal!35!black,title={Before you move on}]
What does mófǎ mean? (\textquotedblleft{}Magic\textquotedblright{}.) Say it once more.
\end{tcolorbox}
\section[bǎowù]{bǎowù — treasure}

\label{lesson:ZH-C331-baowu}



\begin{quote}
Before the new one: say the Mandarin for a Daoist immortal, then the Mandarin for magic in a fairy tale.
\end{quote}

\subsection*{You'll want to know: bǎowù}
\emph{bǎowù} — \textquotedblleft{}treasure\textquotedblright{}.

\begin{quote}
\emph{tā zài shān lǐ zhǎodào le yí ge bǎowù} — He found a treasure in the mountains.
\end{quote}

Its characters come later; for now it is a word for the ear.

\begin{cousinweb}
\emph{bǎowù} is built from \emph{bǎo} 'precious' and \emph{wù} 'a thing', both inherited from Middle Chinese, \emph{pawX} and \emph{mjut} (Guangyun, in Baxter's transcription), which glosses the first 'a treasure' and the second 'all things'.
\end{cousinweb}

\begin{grammarlens}[title={What you've built}]
One more word for telling what happened.
\end{grammarlens}

\subsection*{Your turn}
\begin{itemize}
\raggedright
  \item \emph{Say it:} \emph{bǎowù}
  \item \emph{Say it:} \emph{bǎowù}, once more
  \item \emph{Say it:} \emph{tā zài shān lǐ zhǎodào le yí ge bǎowù} — He found a treasure in the mountains.
  \item \emph{Recall:} say the Mandarin for a Daoist immortal, then the Mandarin for magic in a fairy tale, then say \emph{bǎowù} again
\end{itemize}

\begin{tcolorbox}[breakable,colback=teal!4,colframe=teal!35!black,title={Before you move on}]
What does bǎowù mean? (\textquotedblleft{}Treasure\textquotedblright{}.) Say it once more.
\end{tcolorbox}
\section[cáng]{cáng — to hide}

\label{lesson:ZH-C331-cang}



\begin{quote}
Before the new one: say the Mandarin for magic in a fairy tale, then the Mandarin for a precious treasure.
\end{quote}

\subsection*{You'll want to know: cáng}
\emph{cáng} — \textquotedblleft{}to hide\textquotedblright{}.

\begin{quote}
\emph{yāoguài lái de shíhou tāmen dōu cáng zài mén hòubian} — When the monster came, they all hid behind the door.
\end{quote}

Its characters come later; for now it is a word for the ear.

\begin{cousinweb}
\emph{cáng} is inherited from Middle Chinese \emph{dzang} (Guangyun, in Baxter's transcription), which glosses it 'to conceal' and 'to keep out of sight'. The dictionary also lists \emph{dzangH}, 'a storehouse', a different word.
\end{cousinweb}

\begin{grammarlens}[title={What you've built}]
One more word for telling what happened.
\end{grammarlens}

\subsection*{Your turn}
\begin{itemize}
\raggedright
  \item \emph{Say it:} \emph{cáng}
  \item \emph{Say it:} \emph{cáng}, once more
  \item \emph{Say it:} \emph{yāoguài lái de shíhou tāmen dōu cáng zài mén hòubian} — When the monster came, they all hid behind the door.
  \item \emph{Recall:} say the Mandarin for magic in a fairy tale, then the Mandarin for a precious treasure, then say \emph{cáng} again
\end{itemize}

\begin{tcolorbox}[breakable,colback=teal!4,colframe=teal!35!black,title={Before you move on}]
What does cáng mean? (\textquotedblleft{}To hide\textquotedblright{}.) Say it once more.
\end{tcolorbox}
